\section{Closing the proof-size parameters}\label{sec:main-proof}

Both results of Section~\ref{sec:introduction} combine
Theorem~\ref{thm:clause-transfer} at \(h=3\ell\) with
Theorem~\ref{thm:affine}; they differ only in the choice of \(k\).

\begin{proof}[Proof of Theorem~\ref{thm:main}]
Let \(\ell\ge32\), \(n=2^\ell\), \(m=n+1\), and put
\(x=n/(32768\,\ell^2)\). Suppose a refutation has \(S\le e^x\) nodes.
Theorem~\ref{thm:clause-transfer} with \(h=3\ell\) gives a complete affine
family over \(\mathcal Q_{m,\ell}\) with at most
\(M=3S+\binom{n+1}{2}\) blocks and a PC refutation of degree
\(D=12\ell+1\ge2h+1\). We verify \eqref{eq:affine-room} for
\[
 k=\left\lfloor\frac{n}{32\ell}\right\rfloor .
\]
An induction on \(\ell\ge32\) shows
\(32768\,(2\ell+4)\ell^2\le2^\ell\), that is, \(x\ge2\ell+4\). Hence
\(e^x\ge2^{2\ell+4}=16n^2\). Since \(\binom{n+1}2\le4n^2\),
\[
 4M\le12e^x+16n^2\le13e^x\le e^{2x},\qquad \ln(4M)\le2x .
\]
From \(n\ge64\ell\) we get \(k\ge1\) and \(k+1\le2k\), so
\(n<(k+1)\,32\ell\le64\ell k\). Therefore
\[
 (n+1)\ln(4M)\le2n\cdot2x=\frac{n^2}{8192\,\ell^2}
 \le\frac{n^2}{4096\,\ell^2}\le k^2 .
\]
Moreover \(32\ell k\le n\) gives \(k\le m\), \(4(k-1)<n\), and
\[
 2B=2k(12\ell+2)\le32\ell k\le n .
\]
Theorem~\ref{thm:affine} now says that the family has no PC refutation of
degree \(D\), a contradiction. Hence \(S>e^x\ge2^x\). Both rule conventions
are covered by Theorem~\ref{thm:clause-transfer}.
\end{proof}

\begin{proof}[Proof of Corollary~\ref{cor:main-superpolynomial}]
The corollary follows from Theorem~\ref{thm:main}, since
\(n/(32768\,\ell^2)\ge K\ell\) for large \(\ell\). We also record the
independent argument of version 1, which is the one formalized in the
linked file and uses the smaller \(k\) of Lemma~\ref{lem:parameters}.
Fix \(K>0\) and let \(a=\max\{\lceil K\rceil,2\}\).
Use Corollary~\ref{cor:polylog-exclusion} with parameters \(a,13,1\),
and take \(\ell\) above its uniform threshold and at least two.
Put \(n=2^\ell\). Suppose a refutation has \(S\le n^K\) nodes.
Since \(n\ge1\) and \(a\ge K\), \(S\le n^a\).
Theorem~\ref{thm:clause-transfer}, with \(h=3\ell\), constructs an actual
complete affine registry with
\[
 N_{\mathrm{reg}}\le3S+\binom{n+1}{2},\qquad D=12\ell+1.
\]
Its initial clauses are those of the usual bit CNF, and its empty final
clause gives a PC derivation of one from exactly the complete family over
\(\mathcal Q_{n+1,\ell}\). No additional initial-value hypothesis remains.
Now
\[
 \binom{n+1}{2}\le(n+1)^2\le4n^2\le4n^a,
 \qquad N_{\mathrm{reg}}\le7n^a\le13\,2^{a\ell}+1,
 \qquad D\le13(\ell+1).
\]
The number of proper blocks is at most \(N_{\mathrm{reg}}\), so these are precisely the
inventory and degree bounds excluded by the corollary, a contradiction.
The threshold is independent of the alleged proof. Therefore \(S>n^K\).

This argument uses \(k=\lceil(\ell+1)(\sqrt2)^\ell\rceil\), giving
\begin{equation}\label{eq:main-asymptotics}
 k=O(\sqrt n\log n),\qquad
 k(D+1)=O\bigl(\sqrt n(\log n)^2\bigr)=o(n).
\end{equation}
far below the room \(k\approx n/(32\ell)\) used for Theorem~\ref{thm:main}.
\end{proof}

\subsection{What the conclusion uses}
The source clauses may have any affine rank. The argument partitions
the actual complete block family by rank, uses one common multiplier
on the high side, and packs every block on the low side.
The final contradiction is the impossibility of an old consequence
\(f\ne0\) in a separated row-linear space. It is not an application
of a degree lower bound for one alone.

Only the old functional system uses PC/NS equality, established by
prescribed-moment extension. The augmented system is handled by a
primitive PC replay. Similarly, only the number of source blocks and
their compact input descriptions need be polynomial. Expanded PC
polynomials and their derivations need not have polynomial monomial
count: the contradiction uses degree.

\subsection{Relationship to the broader PHP project}
This result arose within a broader mathematical research project led by
Kamil Braun, whose main objective is a superpolynomial lower
bound for ordinary PHP in every fixed-depth
\(\mathrm{AC}^0[p]\)-Frege system. That objective remains open.
Theorem~\ref{thm:main} concerns a specific subsystem over \(\F\).
The project's research record, supporting code, and autonomous-assistant
workflow are maintained in \emph{Noemesis}, an open framework for autonomous
mathematical research and formal verification. Its source and research archive
are publicly available in the \texttt{math-research} repository on GitHub
\cite{Repo}:
\begin{center}
 \url{https://github.com/kbr-/math-research}
\end{center}
Section~\ref{sec:framework} describes the archive and framework.

\section{A generic sufficient condition}\label{sec:generic}

This section can be skipped. It records that Steps~1 and~2 of the proof
never use the pigeonhole principle, and states what they give for an
arbitrary set of linear clauses. Everything is over \(\F\). Proofs are by
reference to the bit-PHP proofs they abstract and to the formal proofs.
The criterion was first extracted from the argument in an AI review of
version 1 (Claude Fable 5.1, 16 September 2026, recorded in the research
notebook) and was then formalized.

For a linear clause \(A\) with true-indicators \(g_1,\ldots,g_s\), its
\emph{clause polynomial} is \(\prod_i(1-g_i)\), of degree at most the width
\(s\); it vanishes exactly at the assignments satisfying \(A\). For a
polynomial system \(\mathcal F\), call a linear space \(U\) of polynomials
\emph{separated at degree \(B\)} if \(U\cap\Cons_B(\mathcal F)=\{0\}\).

\begin{theorem}[Generic criterion]\label{thm:generic}
\leangroup{\leanref{claims/GenericAffineSubspaceConsequence.lean}{PC level}
  \leansep \leanref{claims/GenericCNFSubspaceCriterion.lean}{proof size}}
Let \(\mathcal I\) be a set of linear clauses in \(v\) variables, and let
\((A_j)_{j\in J}\) be a finite family of linear clauses of width at most
\(w\) such that every clause of \(\mathcal I\) is semantically implied by
some \(A_j\). Let \(\mathcal F\) be a polynomial system containing the
Boolean equations and the clause polynomials of all \(A_j\). Let
\(h,k\ge1\) with \(h(k+1)+1\le v\), put
\[
 D=\max\{2h+w,\,4h+1\},\qquad B=k(D+1),
\]
and let \(U\) be a space of polynomials of degree at most \(k\) that is
separated at degree \(B\). Then every \(\Res\) refutation of \(\mathcal I\)
with \(S\) nodes, under either rule convention of
Section~\ref{sec:introduction}, satisfies
\begin{equation}\label{eq:generic-criterion}
 \dim U\le\Bigl(3S+|J|\Bigr)\binom{v-h(k+1)-1+k}{k}.
\end{equation}
\end{theorem}
\begin{proof}[Proof by reference]
Theorem~\ref{thm:dag-registry} is already stated for arbitrary clause sets
with a supplied covering family; it gives a registry with \(3S+|J|\) slots.
The initial values are derived as in \eqref{eq:initial-value}, with the
clause polynomial of \(A_j\) in place of \(E_{ii'}\), through degree
\(2h+w\). This gives a PC refutation of degree \(D\) from \(\mathcal F\)
and the blocks. Lemma~\ref{lem:hybrid} is already stated for an arbitrary
old system. In place of Lemma~\ref{lem:kernel}, restrict \(U\) to the zero
flat of each proper block of rank above \(h(k+1)\): by
Lemma~\ref{lem:restriction-ideal} each image has dimension at most
\(\binom{v-h(k+1)-1+k}{k}\), so if \eqref{eq:generic-criterion} failed,
some nonzero \(f\in U\) would restrict to zero on all of them, and
Lemma~\ref{lem:hybrid} would give \(f\in\Cons_B(\mathcal F)\).
\end{proof}

\begin{corollary}[Density form]\label{cor:generic-density}
\lean{claims/GenericCNFDensityBound.lean}
Under the hypotheses of Theorem~\ref{thm:generic},
\[
 3S+|J|\ \ge\ \delta\,\exp\!\left(\frac{hk^2}{v+k}\right),
 \qquad \delta=\frac{\dim U}{\binom{v+k}{k}} .
\]
\end{corollary}
\begin{proof}
Removing \(r\) of \(v\) variables shrinks \(\binom{v+k}{k}\) by at least the
factor \(\exp(-rk/(v+k))\), since each of the \(k\) factors
\((v-r+j)/(v+j)\) is at most \(1-r/(v+k)\). Take \(r=h(k+1)+1\ge hk\).
\end{proof}

Here \(\delta\) is the density of \(U\) among all polynomials of degree at
most \(k\). The bound is useful only if \(\delta\) is not much smaller than
\(\exp(-hk^2/(v+k))\) while \(U\) remains separated at degree
\(k(D+1)\). For bit PHP, \(U=\mathcal L_k\), \(w=\ell\), \(h=3\ell\), the old
system is \(\mathcal Q_{m,\ell}\), the covering family is the
\(\binom m2\) compact clauses of Theorem~\ref{thm:clause-transfer}, and
separation is Theorem~\ref{thm:row-separation}; the sharper count of
Lemma~\ref{lem:kernel} then gives Theorem~\ref{thm:main}.
Thus Sections~\ref{sec:affine} and~\ref{sec:clause-simulation} are generic
for any clause set whose covering clause polynomials are axioms of the base.
What is specific to bit PHP is the separation theorem over the compact base,
and the fact that a large separated space exists at all.

\begin{remark}[A sufficient condition, not a tradeoff]\label{rem:not-tradeoff}
Theorem~\ref{thm:generic} is not a relation between \(\Res\) size and PC
refutation degree, nor between size and width. Its hypothesis is not that
\(\mathcal F\) has no low-degree refutation but that a large space \(U\) of
low-degree polynomials meets the low-degree consequences only in zero.
The conclusion of Step~2 is a nonzero derivable member of \(U\), not a
derivation of one, and nonzero polynomials can be legitimate low-degree
consequences: the satisfiable system \(\{x,\ x^2-x\}\) derives \(x\) and
does not derive one. We do not know a formula other than bit PHP for which
a separated space of useful density has been established.
\end{remark}

\begin{remark}[Formulas with short refutations]\label{rem:short-proofs}
\lean{claims/ShortProofControl.lean}
Read contrapositively, a clause set with a size-\(S\) refutation admits no
separated space of degree \(k\) and dimension above the right side of
\eqref{eq:generic-criterion}. As a consistency check on the easiest parity
contradiction: if \(\mathcal F\) contains a polynomial \(L\) of degree at
most one and also \(1-L\), then it derives one at degree one, every
polynomial of degree at most \(B\) is derivable at degree \(B\ge1\), and
the only separated space is zero. This checks one trivial family; it is not
an audit of the known short \(\Res\) refutations.
\end{remark}

\subsection{What does not follow}
\emph{Unary PHP.} The standard substitution from unary PHP into bit PHP
replaces a unary variable by the conjunction describing one bit label; see,
for example, \cite{BB21}. For bounded-depth Frege it has polynomial size
overhead and a constant depth increase. But an affine form in unary
variables becomes a polynomial in bit conjunctions, which is generally
nonlinear, so the substitution does not transfer Theorem~\ref{thm:main} to
unary PHP in \(\Res\). Nor does Theorem~\ref{thm:generic} apply directly:
no sufficiently dense separated space is known for the unary encoding.

\emph{Size and width.} No relation between \(\Res\) size and width, or
between size and PC degree, is claimed (Remark~\ref{rem:not-tradeoff}).

\emph{Stronger systems.} A lower bound for the subsystem \(\Res\) does not
by itself give a Frege lower bound. In the translation of
\(\mathrm{AC}^0[p]\)-Frege proofs \cite{BIKPRS} the extension axioms are
nested: inputs of one block mention variables of others. Step~2 uses that
all inputs are affine in the original variables, and we know of no
removal step for nested blocks.

Other possible applications suggested in review, such as further succinctly
encoded principles or resolution over linear equations modulo a prime
\(p>2\), have not been examined and nothing is claimed about them.

\section{Evidence and the research framework}\label{sec:framework}

The supplementary research record is maintained as a dated online notebook:
\begin{center}
 \url{https://kbr.is-a.dev/math-research/}
\end{center}
It presents arguments, failed attempts, and computation records from the
broader project. The GitHub repository introduced above stores its source,
supporting programs and data, and earlier versions. These materials document
how the result was developed; the mathematical argument is presented in
this manuscript.

\subsection{Mathematical provenance and component checks}
This manuscript presents the
\href{https://kbr.is-a.dev/math-research/\#lean-publication-bit-PHP-superpolynomial}
{formalized notebook theorem and its proof}
\lean{claims/BitPHPSuperpolynomial.lean}, recorded on 15 September 2026,
together with the
\href{https://kbr.is-a.dev/math-research/\#lean-bit-PHP-exponential}
{exponential bound} and the
\href{https://kbr.is-a.dev/math-research/\#lean-generic-CNF-affine-DAG-subspace-criterion}
{generic criterion}, formalized on 17 September 2026.
Version 1 corresponds to repository revision
\href{https://github.com/kbr-/math-research/tree/54f09378e97465522b7f1caf115f90ad1faa7a57}
{\texttt{54f0937}}, which contains the superpolynomial theorem and its complete
formal dependency route. The formal content of revision 1 is contained in
\href{https://github.com/kbr-/math-research/tree/b47e9b1ef1f5b273822001983e83d35d7acbe117}
{\texttt{b47e9b1}}, which adds the modules for Theorem~\ref{thm:main},
Theorem~\ref{thm:generic}, Corollary~\ref{cor:generic-density}, and
Remark~\ref{rem:short-proofs}.
The earlier
\href{https://kbr.is-a.dev/math-research/\#bit-PHP-publication-audit-verdict}
{internal publication audit} and subsequent literature notes in the
notebook distinguish the
claimed unrestricted result from the regular and depth-restricted
literature.

The repository retains exact primitive-PC traces for the rule
simulation, including missing-premise controls.
A separate matching-moment checker stores sparse separating
functionals and every potentially nonzero marginal equation for
finite instances, together with falsifying perturbations.
These finite experiments were development checks, not a proof of the asymptotic
statement. The subsequent Lean formalization now verifies that statement
and its entire dependency route, including the chessboard filling theorem,
actual CNF-to-PC translation, and uniform parameter closure. The permanent
[Lean] links in this revision refer to published commit
\href{https://github.com/kbr-/math-research/tree/b47e9b1ef1f5b273822001983e83d35d7acbe117}
{\texttt{b47e9b1}}.
Its project pins Lean \texttt{4.34.0-rc2} and the Mathlib dependency revision.
The final transitive axiom report contains only Lean's standard foundations;
there is no \texttt{sorry}, custom axiom, or unproved chessboard interface.
The compiled module of the superpolynomial theorem and all its imported
declarations also passed a fresh kernel replay with the bundled
\texttt{leanchecker --fresh} on 15 September 2026; the modules added in
revision 1 passed the project's verification script, with the same axiom
report, but were not separately replayed in this way. The complete project was separately rebuilt from empty
project build outputs using the pinned library caches.
The proofs here follow the formalization entries, including
alternative arguments adopted during formalization. Formal verification of
the stated objects is distinct from external assessment of the manuscript's
exposition, interpretation of the proof system, and novelty.

\subsection{An inspectable workflow for autonomous research}
The project also develops \emph{Noemesis}, an open framework for autonomous
mathematical research: a collection of instructions, tools, and record-keeping conventions
for AI assistants, distributed in the same repository.
Its purpose is to let an autonomous assistant continue a mathematical
investigation across finite context windows and session boundaries
while retaining inspectable evidence \cite{Repo}.
Its main mechanisms are:
\begin{itemize}
 \item \textbf{Persistent context.}
 A stable restart guide, a concise living overview, and a claim index
 recover the relevant definitions and proofs without reloading the
 full conversation or depending on a particular machine.
 \item \textbf{Recorded successes and failures.}
 Dated, append-only entries retain arguments, assumptions,
 obstructions, and unsuccessful attempts with explicit claim status.
 Corrections receive new entries instead of silently replacing history.
 \item \textbf{Reproducible experiments.}
 Protected execution, complete outputs, provenance hashes, and
 timing records preserve what was actually computed under bounded
 resources.
 \item \textbf{Repeated investigation and process review.}
 Each cycle targets a named obligation, checks its outcome, and
 assesses wasted work or missing tools. Bounded framework
 improvements are recorded separately from mathematical results.
 \item \textbf{Explicit formal verification.}
 Assigned formalization cycles build dependency-indexed Lean proofs, record
 their exact scope and arguments, and link them from the
 claim index. Independent branches can run in isolated worktrees before
 integration; cached dependencies avoid redundant verification.
 \item \textbf{A live, versioned notebook.}
 Saved versions in Git, a version-control system, and the public notebook let
 readers inspect the development locally or online.
\end{itemize}
The durable record uses ordinary source, HTML, Markdown, and data files.
These mechanisms make work recoverable and reviewable; they do not
make a recorded proof automatically correct or formally verified.

\subsection{Contributions and AI assistance}\label{sec:model-assistance}
The mathematics in this paper was produced by AI models working inside the
framework above, under my direction. For version 1, GPT-6 Astra (in the
Codex command-line agent) was the main contributor to the mathematical
development, the Lean formalization, and the drafting; Claude Sonnet 5,
Claude Fable 5.1, and Claude Opus 5 were used in earlier exploratory
conversations. For revision 1, Claude Fable 5.1 (in Claude Code) formalized
the exponential bound of Theorem~\ref{thm:main} and the generic criterion of
Section~\ref{sec:generic}, statements and proofs, with one batch of proofs
delegated to Claude Opus 5; it also triaged reader and AI-review feedback on
version 1 and drafted the revised text, including the proof overview and
the related-work discussion. My own contribution is the design, direction,
and iterative refinement of the research framework, the choice of problems,
and steering questions about whether the work was approaching its goal; it
is not the mathematical derivations. I do not have the expertise to check
the mathematics myself, and I did not check it, in version 1 or in this
revision. My confidence in the results rests on the Lean formalization,
whose scope is the statements carrying [Lean] links, under the definitions
in the linked files. The prose, the overview, the statements about the
literature, and the fidelity of the formal definitions to the standard
proof system are not covered by it. Neither version has been externally
peer reviewed.

A first-person account of how the project developed (the chat
conversations it began in, the move to file-based agents, the
\emph{Spin} research loop, and the parallel formalization) appeared as
Section~8.3 of version 1. It has been moved out of the paper and is
preserved in the repository as
\href{https://github.com/kbr-/math-research/blob/main/publications/bit-php-resolution-over-parities/DEVELOPMENT_HISTORY.md}
{\texttt{DEVELOPMENT\_HISTORY.md}}; the version-1 source containing it is
fixed at commit
\href{https://github.com/kbr-/math-research/blob/b47e9b1ef1f5b273822001983e83d35d7acbe117/publications/bit-php-resolution-over-parities/sections/07-main-proof-framework.tex}
{\texttt{b47e9b1}}.

The classical topological input is attributed to \cite{BLVZ94};
the product-extension methodology is attributed to \cite{BIKPRS}.
Original manuscript material is distributed under the
\href{https://creativecommons.org/licenses/by/4.0/}
{Creative Commons Attribution 4.0 International license (CC BY 4.0)}.
